\documentclass[aps,prb,twocolumn,superscriptaddress,floatfix]{revtex4-2}
\usepackage{amsmath,amssymb}
\newcommand{\rp}{r_\perp}
\renewcommand{\thetable}{S\arabic{table}}
\renewcommand{\thesection}{S\arabic{section}}
\begin{document}
\title{Supplemental Material for ``Onset of a quantum Fisher--Selke sequence in the transverse-field ANNNI model and its
bilayer''}
\maketitle

\section{Exact-diagonalization benchmark}

Table~\ref{tab:ed} compares the SSE code with exact diagonalization at finite temperature on four 12-spin periodic
clusters built by the production lattice code (Sec.~IV~C of the main text). The benchmark anneals every case over 16 fields from $g=3$ to the target,
with 2000 sweeps per field, followed by 5000 thermalization and 40\,000 measurement sweeps at the
target, in 24 independent chains.

\begin{table*}
\caption{SSE against exact diagonalization on 12-spin periodic clusters. $e$ is the energy per spin
and $m_x$ the transverse magnetization, both in units of $J_0$; the SSE values are means of 24
chains with standard errors in parentheses. The last column is the largest $|z|$ over all
observables of the case (31--39 per case).}
\label{tab:ed}
\begin{ruledtabular}
\begin{tabular}{llcrrrrlllll}
Cluster & & $T/J_0$ & $J_y/J_0$ & $\kappa$ & $\rp$ & $g$ & $e$ (ED) & $e$ (SSE) & $m_x$ (ED) & $m_x$ (SSE) & max$|z|$ \\
\hline
$2\times3$ bilayer & & 0.15 & 1 & 0 & 1 & 1.0 & $-2.60042$ & $-2.60075(20)$ & 0.20169 & 0.20176(14) & 1.65 \\
$2\times3$ bilayer & & 0.15 & 1 & 0.75 & $-1$ & 0.5 & $-1.78588$ & $-1.78603(17)$ & 0.14416 & 0.14394(23) & 2.01 \\
$2\times3$ bilayer & & 0.50 & 1 & 0.5 & 1 & 2.0 & $-2.52091$ & $-2.52076(68)$ & 0.54569 & 0.54629(57) & 2.62 \\
$2\times6$ layer & & 0.15 & 1 & 0 & 0 & 2.0 & $-2.52595$ & $-2.52599(23)$ & 0.56144 & 0.56135(18) & 0.68 \\
$2\times6$ layer & & 0.15 & 1 & 0.75 & 0 & 1.0 & $-1.83990$ & $-1.84065(41)$ & 0.46098 & 0.46162(69) & 1.83 \\
$2\times6$ layer & & 0.15 & 1 & 0.5625 & 0 & 1.5 & $-2.07157$ & $-2.07171(41)$ & 0.64328 & 0.64283(47) & 1.86 \\
$2\times6$ layer & & 0.15 & $-1$ & 0.25 & 0 & 1.0 & $-1.89600$ & $-1.89616(20)$ & 0.29902 & 0.29881(20) & 1.82 \\
$2\times6$ layer & & 0.50 & 1 & 1.0 & 0 & 0.5 & $-1.71530$ & $-1.71590(69)$ & 0.18847 & 0.18876(157) & 2.00 \\
$1\times6$ ladder & & 0.15 & 1 & 0.75 & 1 & 1.0 & $-1.46156$ & $-1.46128(36)$ & 0.61156 & 0.61140(57) & 1.80 \\
$1\times6$ ladder & & 0.15 & 1 & 1.3 & $-1$ & 0.3 & $-1.30261$ & $-1.30361(86)$ & 0.22459 & 0.22781(527) & 2.50 \\
$1\times6$ ladder & & 0.15 & 1 & 0.5 & 1 & 3.0 & $-3.13579$ & $-3.13617(44)$ & 0.95525 & 0.95528(12) & 1.68 \\
$3\times4$ layer & & 0.15 & 1 & 0.375 & 0 & 2.5 & $-2.74079$ & $-2.74034(26)$ & 0.88540 & 0.88520(14) & 1.68 \\
$3\times4$ layer & & 0.15 & 1 & 1.0 & 0 & 1.0 & $-2.12664$ & $-2.12642(20)$ & 0.25671 & 0.25629(20) & 2.08 \\
\end{tabular}
\end{ruledtabular}
\end{table*}

\section{Pre-registered test of the $\langle23\rangle$ phase}

The test of Sec.~V~B of the main text followed a written plan. Its design, analysis and pass criteria were fixed before
the production runs, after an independent review of the design.

\paragraph*{Design.} Seven points $(\kappa,r)$ with $\lambda\le0.5$, chosen so that the predicted margin
of $\langle23\rangle$ at $g_c$ exceeds the expected statistical error. At each point the pairs
$\{\langle23\rangle,\langle2\rangle\}$ ran on a $20\times20$ bilayer and $\{\langle23\rangle,\langle3\rangle\}$ on
a $20\times30$ bilayer, with 12 chains per structure, $2000+3000$ sweeps per field at $T=0.15\,J_0$, and
the ascending template protocol of Sec.~IV~B of the main text. The field grid had steps of $0.1$ below the
predicted window, steps of $0.02$ (or $0.01$) across it with a margin of at least $0.06$, and $g_c$ as an
exact grid point. Competitors ($\langle223\rangle$, $\langle233\rangle$, $\langle2333\rangle$) ran at four of
the points, and one point was repeated at $T=0.10\,J_0$.

\paragraph*{Analysis.} Free energies as in Sec.~IV~B of the main text. Uncertainties from the jackknife
over chains; differences tested with Welch's $t$ and its Satterthwaite degrees of freedom; crossings
from a cubic-spline root of the free-energy difference, with linear interpolation in $g^2$ as a cross
check.

\paragraph*{Criteria.}
(T1)~Existence at each point: $F_{\langle23\rangle}<F_{\langle2\rangle}$ and $F_{\langle23\rangle}<F_{\langle3\rangle}$
at $g_c$, each with one-sided $p<0.0228$, and no chain cut before $g_c$. The headline required a pass at
all seven points, and a failure at $\lambda\le0.4$ was to be reported and investigated.
(T2)~Window edges compared with the exact-$\kappa$ $O(g^4)$ and $O(g^6)$ values (not pass/fail).
(T3, T4)~Width of the window and the ratio of widths for $r=0$ and $1$.
(T5)~Competitors at $g_c$ against their $O(g^6)$ values, with a stop rule: no claim for
$\langle23\rangle$ if a competitor lay below it at $g_c$ by more than $3\sigma$.
(T6)~Consistency of $\langle23\rangle$ on the two lattices ($\max|z|<3$) and of $\langle3\rangle$ with the
$24\times24$ branches.
(T7)~Window edges at $T=0.10$ and $0.15$ agree within $2\sigma$.
(T8)~An exploratory search for further structures above the window.

Table~\ref{tab:prereg} lists the outcomes. The stop rule of T5 was never triggered. At $\kappa=0.53$, $r=1$, T6
passed its criterion, but with $\chi^2=3.45$ per point and 14 of 25 fields beyond $2\sigma$; the original $20\times30$ run
for $\langle23\rangle$ at this point is the one that the repeat with four times the sweeps did not reproduce
(Sec.~V~B of the main text).

\begin{table}
\caption{Outcomes of the pre-registered criteria. T2: pulls of the lower and upper edges against $O(g^6)$. T5: largest
$|{\rm pull}|$ of a competitor against $O(g^6)$. T6: $\max|z|$ of $\langle23\rangle$ on $20\times20$ against
$20\times30$, with $\chi^2$ per point. T7: pulls of the edges at $T=0.10\,J_0$ against $T=0.15\,J_0$. T8: fields at which
$\langle233\rangle$ lies $2\sigma$ below $\langle23\rangle$ and $\langle3\rangle$. T3 and T4 are in
Table~VI of the main text and Fig.~5(c) of the main text.}
\label{tab:prereg}
\footnotesize
\begin{ruledtabular}
\begin{tabular}{llllllll}
$\kappa$ & $r$ & T1 & T2 & T5 & T6 & T7 & T8 \\
\hline
0.52   & 0 & pass & $-1.8$, $-1.1$  & ---   & 1.15 (0.57) & --- & --- \\
0.53   & 0 & pass & $+0.1$, $+2.7$  & 1.4   & 1.35 (0.59) & $-0.1$, $-1.7$ & none \\
0.545  & 0 & pass & $-0.7$, $+5.1$  & 1.2   & 1.85 (0.55) & --- & none \\
0.5625 & 0 & pass & $-0.9$, $+11.4$ & 3.3   & 1.04 (0.34) & --- & 1.10--1.30 \\
0.52   & 1 & fail & $+1.4$, $+1.2$  & ---   & 1.38 (1.14) & --- & --- \\
0.53   & 1 & pass & $-0.5$, $+2.9$  & 1.5   & 2.19 (3.45) & --- & none \\
0.545  & 1 & pass & $+0.4$, $+1.4$  & ---   & 1.15 (0.34) & --- & --- \\
\end{tabular}
\end{ruledtabular}
\end{table}

\paragraph*{Changes.} Before the runs, the review moved $\langle3\rangle$ from a $12\times12$ to the
$20\times30$ lattice, added $g_c$ to the grids, specified the quadrature, and replaced a test on
$\min(F_{\langle2\rangle},F_{\langle3\rangle})$ by the two separate comparisons of T1. After the first
output existed, eight ambiguities in the analysis were resolved in writing without changing the text of
the plan, and an error that affected only chains cut at a change of structure was corrected. It had no
effect on any chain used in T1. The follow-up at $\kappa=0.52$, $r=1$ was decided after the result of
T1 at that point was known, and is reported separately.

\section{Additional runs and error budget}

\subsection{Common lattice, replicates and sweeps}
\label{sec:s3runs}

After the pre-registered test (Sec.~V~B of the main text) we ran the following checks; Table~\ref{tab:lattices}
collects the results. For decoupled layers the common lattice at $\kappa=0.53$ places the upper edge $1.7\sigma$ above
the eighth-order value and $2.9\sigma$ above the sixth. In the bilayer the pre-registered upper edges lie above the
series, and the paired differences of $e_b-f_b$ between $\langle23\rangle$ and $\langle3\rangle$ are consistent with
zero, so the integration does not cause this excess.

To locate the bilayer excess we followed all three structures on one $20\times60$ lattice at both bilayer points, twice
with independent seeds, and repeated $\langle23\rangle$ and $\langle3\rangle$ with four times the sweeps on both lattices
at $\kappa=0.52$ (Table~\ref{tab:lattices}). For decoupled layers the common lattice agrees with the pre-registered
values and with the eighth-order series at both $\kappa$; at $\kappa=0.52$ it also turns the narrowest pre-registered
margin, $\langle23\rangle$ below $\langle3\rangle$, into a clear one. In the bilayer the two $20\times60$ runs agree with
each other, and together they place both edges and both gains within $2\sigma$ of the eighth-order series. The gain over
$\langle2\rangle$ is the same on every lattice. The gain over $\langle3\rangle$ is about $10^{-4}$ per spin larger on
$20\times30$ at the standard sweeps, in every set run with these settings: the original runs, the follow-up at
$\kappa=0.52$ and the repeat at $T=0.10\,J_0$. With four times the sweeps it falls to the $20\times60$ value at both
points, while on $20\times60$ it does not change. At $\kappa=0.52$ the four combinations of lattice and sweeps thus
place the shift in the one combination of $20\times30$ and standard sweeps (interaction $1.9\sigma$, simple effects
$2.1$--$2.5\sigma$), and we think these runs were not fully equilibrated. The original $20\times30$ run for
$\langle23\rangle$ at $\kappa=0.53$ had also been flagged by the consistency check T6 (Appendix~C of the main text). The
energies shift with the free energies, so the integration is not the cause, and the pair correlations along $y$ show no
wall defects. Runs repeated with identical settings and new seeds agree within their statistical errors (rms $t=1.06$
over ten comparisons). We therefore take the $20\times60$ and four-times-sweeps values as our estimates. They agree with
the eighth-order series within $1\sigma$ in the gain over $\langle3\rangle$ and the upper edge at both points; in the
bilayer the sixth and eighth orders differ by less than these errors. The pre-registered $\Delta F_3$ in the bilayer
carries a bias of about $-10^{-4}$ per spin, which raises the upper edge by about $0.008$ in $g$. $\langle23\rangle$ lies
below both neighbors in every run and on every lattice at all seven points. At $\kappa=0.52$, $r=1$ its margin over
$\langle3\rangle$ is $1.6$--$1.9\times10^{-4}$ per spin in the three unbiased sets, five to nine standard errors, and
its margin over $\langle2\rangle$ on the common lattice is $1.6\times10^{-4}$, $6.5$ standard errors.

At $T=0.10\,J_0$, where the thermal contributions estimated below are about ten times smaller, $\langle23\rangle$
remains below both neighbors at $\kappa=0.53$ for both $r$ (Table~\ref{tab:lattices}). For $r=0$ the lower edge and the
gain over $\langle2\rangle$ are unchanged, and the upper edge and the gain over $\langle3\rangle$ move by $1.7\sigma$ and
$1.5\sigma$; for $r=1$ both edges and both gains agree with those at $T=0.15\,J_0$ within $1\sigma$.

\begin{table}
\caption{Test points on different lattices and in different runs: window edges and the free-energy
differences $\Delta F_2$ and $\Delta F_3$ of $\langle23\rangle$ relative to $\langle2\rangle$ and $\langle3\rangle$ at $g_c$
(in $10^{-4}J_0$ per spin, negative where $\langle23\rangle$ is lower), with statistical errors. The
pre-registered runs use $20\times20$ for $\langle2\rangle$ and $20\times30$ for $\langle3\rangle$; the follow-up at
$\kappa=0.52$ was decided afterwards; ``$4\times$'' has four times the sweeps; $20\times60$ holds all three
structures, ``replicate'' repeats that run with new seeds, and ``pooled'' combines the chains of both; ``$T=0.10$'' repeats the
pre-registered runs at $T=0.10\,J_0$; $40\times40$ holds $\langle2\rangle$ and $\langle23\rangle$ only.}
\label{tab:lattices}
\footnotesize\setlength{\tabcolsep}{1.2pt}
\begin{ruledtabular}
\begin{tabular}{llllll}
$\kappa$, $r$ & run & $g_{2|23}$ & $g_{23|3}$ & $\Delta F_2$ & $\Delta F_3$ \\
\hline
0.52, 0 & pre-registered & 0.5885(21) & 0.6084(36) & $-3.32(60)$ & $-1.46(62)$ \\
     & $20\times60$ & 0.5917(13) & 0.6130(28) & $-2.34(38)$ & $-2.26(47)$ \\
     & $O(g^8)$ & 0.5922 & 0.6134 & $-2.21$ & $-2.32$ \\
0.53, 0 & pre-registered & 0.7050(19) & 0.7553(51) & $-4.93(76)$ & $-7.19(99)$ \\
     & $20\times60$ & 0.7038(11) & 0.7522(37) & $-5.41(42)$ & $-6.84(72)$ \\
     & $T=0.10$ & 0.7048(17) & 0.7451(33) & $-5.12(65)$ & $-5.36(66)$ \\
     & $O(g^8)$ & 0.7045 & 0.7460 & $-5.16$ & $-5.63$ \\
0.53, 1 & pre-registered & 1.1037(29) & 1.1604(27) & $-4.31(66)$ & $-4.90(34)$ \\
     & $20\times30$, $4\times$ & --- & 1.1550(29) & --- & $-4.23(37)$ \\
     & $20\times60$ & 1.1033(19) & 1.1483(39) & $-4.38(45)$ & $-3.38(51)$ \\
     & replicate & 1.1017(18) & 1.1552(32) & $-4.75(40)$ & $-4.25(43)$ \\
     & pooled & 1.1025(13) & 1.1517(26) & $-4.56(30)$ & $-3.82(34)$ \\
     & $40\times40$ & 1.1049(11) & --- & $-4.03(27)$ & --- \\
     & $T=0.10$ & 1.1069(17) & 1.1612(37) & $-3.56(39)$ & $-5.06(49)$ \\
     & $O(g^8)$ & 1.1046 & 1.1536 & $-4.08$ & $-4.09$ \\
0.52, 1 & pre-registered & 0.9297(30) & 0.9577(52) & $-0.99(54)$ & $-2.47(58)$ \\
     & follow-up & 0.9232(25) & 0.9620(36) & $-2.13(46)$ & $-2.90(39)$ \\
     & $20\times60$ & 0.9258(19) & 0.9477(40) & $-1.67(34)$ & $-1.38(44)$ \\
     & replicate & 0.9269(20) & 0.9527(38) & $-1.48(35)$ & $-1.90(41)$ \\
     & pooled & 0.9264(13) & 0.9502(28) & $-1.58(24)$ & $-1.64(30)$ \\
     & $20\times30$, $4\times$ & --- & 0.9526(23) & --- & $-1.91(25)$ \\
     & $20\times60$, $4\times$ & --- & 0.9524(20) & --- & $-1.87(22)$ \\
     & $O(g^8)$ & 0.9253 & 0.9515 & $-1.78$ & $-1.78$ \\
\end{tabular}
\end{ruledtabular}
\end{table}

\subsection{Error budget}

The error budget of the free-energy differences at the test points has five parts. The statistical errors are
$0.3$--$1.2\times10^{-4}$ per spin (Table~VI of the main text). The choice of quadrature (trapezoidal, Richardson or spline)
moves them by up to $0.4\times10^{-4}$. The truncation of the series, estimated by the last retained term, the
difference between the eighth- and sixth-order values of Tables~VI and VII of the main text, is at most
$0.5\times10^{-4}$ in the bilayer and $0.7\times10^{-4}$ for $r=0$ at $\kappa\le0.53$, but $2.6\times10^{-4}$ and
$7.6\times10^{-4}$ for $\Delta F_3$ at $\kappa=0.545$ and $0.5625$, $r=0$. Thermal excitations of the rows next to the
walls contribute at $T=0.15\,J_0$. For $r=0$ a free-chain estimate, which neglects the pinning of these rows by the walls
and therefore bounds the effect from above, gives up to $0.3\times10^{-4}$ at $\kappa=0.52$, $1.0\times10^{-4}$ at
$\kappa=0.53$ and several $10^{-4}$ at $\kappa\ge0.545$; it favors $\langle23\rangle$ over $\langle2\rangle$ and
$\langle3\rangle$ over $\langle23\rangle$, the opposite of the upper-edge excess. For $r=1$ the rows form ladders whose kink
energy at the test fields exceeds $1.2J_0$, and the estimate is below $10^{-5}$. Finally, the pre-registered $20\times30$ runs in the bilayer carry the
bias of about $-10^{-4}$ in $\Delta F_3$ discussed in Sec.~\ref{sec:s3runs}. For $r=0$ at $\kappa\ge0.545$ the comparison of $\Delta F_3$ and of the upper
edge with the series is therefore qualitative.

\end{document}
